Cooperative multi-agent reinforcement learning under partial observability is
typically trained under the centralised-training / decentralised-execution
(CTDE) regime \citep{rashid2018qmix, yu2022mappo}. Two attention-based ideas
have dominated the design space. The first is the centralised \emph{value
mixer} \citep{yang2020qatten, rashid2018qmix}: a non-linear, often
attention-driven combination of per-agent utilities into a joint
$Q_{\text{tot}}$ that preserves a tractable decentralised argmax. The second
is the \emph{communication channel} \citep{das2019tarmac, jiang2018atoc,
sukhbaatar2016commnet, foerster2016dial}: per-agent learned messages routed
through soft attention so that each agent's policy can condition on a
state-dependent summary of peer information at every step.

The communication channel is a powerful primitive — TarMAC
\citep{das2019tarmac} reports clear gains over broadcast and over no-comm
baselines on traffic-junction and predator-prey tasks — but every realisation
we reviewed is dense: each agent computes attention weights and aggregates
messages from all $N{-}1$ peers, an $\mathcal{O}(N^2)$ cost per simulation
step. For distributed MARL infrastructure
\citep{liang2017rllib, espeholt2020seedrl}, where the per-step inference
dominates total wall-clock, this cost is the binding constraint at large
$N$.

This paper asks a sharply scoped question: \emph{can the soft attention
channel be sparsified to top-$k$ destinations without sacrificing return,
and can $k$ itself be learned per agent per step?} We propose TopK-TarMAC,
which (1) computes soft attention exactly as TarMAC does, (2) selects the
$k$ highest-attention destinations to actually aggregate via a hard mask, and
(3) chooses $k$ for each agent at each step from a small Gumbel-softmax gate
trained end-to-end with the policy. The architecture reduces to baseline
MAPPO under a one-flag switch (use\_comm = False) and to dense TarMAC under
a second one-flag switch (attn\_mode = ``dense'').

Our contributions are:
\begin{itemize}
\item TopK-TarMAC, an adaptive-sparsity attention-communication module
  that interpolates between MAPPO ($k{=}0$ messages) and dense TarMAC
  ($k{=}N{-}1$), with the gate trained end-to-end via Gumbel-softmax.
\item A 5-seed comparative study at $N \in \{3, 6\}$ and a 3-seed study at
  $N{=}12$, with per-seed final returns, Welch's $t$-test $p$-values,
  Cohen's $d$, and 10\,000-resample percentile bootstrap CIs on the mean
  delta, matching the rigour standard of recent MARL benchmarks
  \citep{papoudakis2021epymarl}.
\item Empirical evidence that attention communication is a real lever at
  moderate $N$ but not at small $N$: dense Attn-Comm yields $d{=}1.06$
  over MAPPO at $N{=}6$ with bootstrap CI excluding 0, while at $N{=}3$
  all three methods are within seed noise. We report this honestly
  rather than tune away the null result.
\item A direct accounting of the FLOPs / return trade-off: the
  adaptive-$k$ gate empirically selects $k_{\mathrm{eff}}/(N{-}1)
  \approx 0.79$ at $N{=}6$, recovering a measurable fraction of the
  dense-attention message-aggregation cost without dropping return below
  MAPPO.
\end{itemize}
